\documentclass[a4paper,11pt]{amsart}

% =========================================================================
%  Informe de laboratorio — Práctica 8: Efecto Doppler con ultrasonidos
%  Estilo: artículo científico minimalista (tipografía Palatino, versalitas,
%  enunciados numerados por sección), paleta de azules y grises sobrios.
% =========================================================================

% --- IDIOMA Y CODIFICACIÓN ------------------------------------------------
\usepackage[utf8]{inputenc}
\usepackage[T1]{fontenc}
\usepackage[spanish,es-tabla,es-nodecimaldot]{babel}

% --- TIPOGRAFÍA (Palatino, como el documento de referencia) ---------------
\usepackage{mathpazo}
\linespread{1.05}
\usepackage[final]{microtype}

% --- MATEMÁTICAS Y UNIDADES ------------------------------------------------
\usepackage{amsmath, amssymb, mathtools}
\usepackage{bm}
\usepackage{siunitx}
\sisetup{
  output-decimal-marker = {,},
  group-minimum-digits  = 4,
  group-separator       = {\,},
  per-mode              = symbol,
  separate-uncertainty  = true,
  range-phrase          = {\text{ a }},
  detect-all            = false
}

% --- PÁGINA, COLORES Y GRÁFICOS --------------------------------------------
\usepackage[a4paper,margin=2.6cm]{geometry}
\usepackage{graphicx}
\usepackage{xcolor}
\definecolor{azulprof}{HTML}{1B365D}    % azul marino
\definecolor{azulacero}{HTML}{3B6EA8}   % azul acero
\definecolor{azulclaro}{HTML}{B7CBE3}   % azul hielo
\definecolor{grispizarra}{HTML}{5B6770} % gris pizarra
\definecolor{grisclaro}{HTML}{E4E8EC}   % gris perla
\usepackage{tikz}
\usetikzlibrary{babel, arrows.meta, calc, positioning, decorations.pathmorphing,
                decorations.markings, patterns}
\usepackage{pgfplots}
\pgfplotsset{compat=1.17}
\usepackage{booktabs}
\usepackage{array}
\usepackage{subcaption}
\usepackage{enumitem}
\setlist[itemize]{leftmargin=1.6em, itemsep=2pt, topsep=4pt}
\setlist[enumerate]{leftmargin=1.8em, itemsep=2pt, topsep=4pt}

% --- ENLACES Y REFERENCIAS CRUZADAS ----------------------------------------
\usepackage[colorlinks=true, linkcolor=azulacero, citecolor=azulacero,
            urlcolor=azulacero, plainpages=false, pdfusetitle=false,
            pdftitle={Práctica 8: Efecto Doppler con ultrasonidos},
            pdfauthor={David Gómez García; Daniel Martín Guerrero}]{hyperref}
\usepackage[spanish,nameinlink]{cleveref}

% --- BIBLIOGRAFÍA -----------------------------------------------------------
\usepackage{csquotes}
\usepackage[backend=biber, style=numeric, sorting=none, giveninits=true,
            maxbibnames=6, minbibnames=3, doi=true, url=true,
            isbn=false]{biblatex}
\addbibresource{referencias.bib}
\DeclareFieldFormat[article,techreport,unpublished]{title}{#1\isdot}
\DeclareFieldFormat{url}{\url{#1}}
\renewcommand*{\bibfont}{\small}
\setlength{\bibitemsep}{3pt}

% --- CABECERAS ---------------------------------------------------------------
\pagestyle{headings}
\makeatletter
\def\@evenhead{\small\scshape\thepage\hfil Laboratorio de Mecánica y Ondas\hfil}
\def\@oddhead{\small\scshape\hfil Práctica 8: Efecto Doppler con ultrasonidos\hfil\thepage}
\def\ps@headings{%
  \let\@mkboth\markboth
  \def\@evenhead{\small\scshape\normalfont\thepage\hfil Laboratorio de Mecánica y Ondas}%
  \def\@oddhead{\small\scshape\normalfont Efecto Doppler con ultrasonidos\hfil\thepage}%
  \let\@evenfoot\@empty \let\@oddfoot\@empty}
\makeatother
\numberwithin{equation}{section}
\setcounter{tocdepth}{1}
\makeatletter
\def\@@and{y}
\renewcommand{\keywordsname}{Palabras clave}
\makeatother

% =========================================================================
%               NOTACIÓN Y MACROS DE USO FRECUENTE
% =========================================================================
\newcommand{\vs}{v_{\mathrm{s}}}              % velocidad del sonido
\newcommand{\vst}{\tilde{v}_{\mathrm{s}}}     % valor teórico esperado
\newcommand{\vE}{v_{\mathrm{E}}}              % velocidad del emisor
\newcommand{\vR}{v_{\mathrm{R}}}              % velocidad del receptor
\newcommand{\fE}{f_{\mathrm{E}}}              % frecuencia emitida
\newcommand{\fR}{f_{\mathrm{R}}}              % frecuencia recibida
\newcommand{\TE}{T_{\mathrm{E}}}
\newcommand{\TR}{T_{\mathrm{R}}}
\newcommand{\dif}{\mathrm{d}}
\newcommand{\Tcoef}{t_{95,\nu}}               % factor de cobertura al 95 %
\newcommand{\ext}[1]{U(#1)}                   % incertidumbre ampliada
\newcommand{\mean}[1]{\overline{#1}}
\DeclareMathOperator{\var}{var}
\DeclareMathOperator{\cov}{cov}
% Dato pendiente de incorporar (se resalta con discreción)
\newcommand{\pend}[1]{\textcolor{grispizarra}{\colorbox{grisclaro}{\strut\,#1\,}}}
\newcommand{\pendiente}[1]{\textcolor{grispizarra}{\textit{[#1]}}}

% =========================================================================
%        ENTORNOS FORMALES (enunciados numerados, estilo artículo)
% =========================================================================
\usepackage{amsthm}

% Enunciados en cursiva (resultados generales)
\newtheoremstyle{cienplano}{10pt}{10pt}{\itshape}{}%
  {\bfseries\color{azulprof}}{.}{0.6em}%
  {\thmname{#1}\thmnumber{ #2}\thmnote{\normalfont\upshape\ (#3)}}
% Definiciones, cálculos y estimaciones (cuerpo en redonda)
\newtheoremstyle{ciendef}{10pt}{10pt}{\normalfont}{}%
  {\bfseries\color{azulprof}}{.}{0.6em}%
  {\thmname{#1}\thmnumber{ #2}\thmnote{\normalfont\upshape\ (#3)}}
% Observaciones (encabezado en cursiva)
\newtheoremstyle{cienobs}{8pt}{8pt}{\normalfont}{}%
  {\itshape\color{azulprof}}{.}{0.6em}%
  {\thmname{#1}\thmnumber{ #2}\thmnote{\normalfont\upshape\ (#3)}}

\theoremstyle{cienplano}
\newtheorem{proposicion}{Proposición}[section]
\newtheorem{resultado}[proposicion]{Resultado}
\newtheorem{conclusion}[proposicion]{Conclusión}
\theoremstyle{ciendef}
\newtheorem{definicion}[proposicion]{Definición}
\newtheorem{hipotesis}[proposicion]{Hipótesis}
\newtheorem{calculo}[proposicion]{Cálculo de incertidumbre}
\newtheorem{estimacion}[proposicion]{Estimación lineal}
\theoremstyle{cienobs}
\newtheorem{observacion}[proposicion]{Observación}

\crefname{proposicion}{proposición}{proposiciones}
\Crefname{proposicion}{Proposición}{Proposiciones}
\crefname{resultado}{resultado}{resultados}
\crefname{conclusion}{conclusión}{conclusiones}
\crefname{definicion}{definición}{definiciones}
\crefname{hipotesis}{hipótesis}{hipótesis}
\crefname{calculo}{cálculo}{cálculos}
\crefname{estimacion}{estimación}{estimaciones}
\crefname{observacion}{observación}{observaciones}
\crefname{section}{sección}{secciones}
\Crefname{section}{Sección}{Secciones}
\crefname{subsection}{sección}{secciones}
\Crefname{subsection}{Sección}{Secciones}
\crefname{table}{tabla}{tablas}
\Crefname{table}{Tabla}{Tablas}
\crefname{figure}{figura}{figuras}
\Crefname{figure}{Figura}{Figuras}
\crefname{equation}{ecuación}{ecuaciones}
\Crefname{equation}{Ecuación}{Ecuaciones}
\crefformat{equation}{ecuación~(#2#1#3)}
\Crefformat{equation}{Ecuación~(#2#1#3)}
\crefrangeformat{equation}{ecuaciones~(#3#1#4)--(#5#2#6)}
\crefmultiformat{equation}{ecuaciones~(#2#1#3)}{ y~(#2#1#3)}{, (#2#1#3)}{ y~(#2#1#3)}

% Pies de figura y tabla, en la línea del documento de referencia
\usepackage[font=small, labelfont={sc,color=azulprof}, labelsep=period,
            margin=0.6cm, justification=justified]{caption}
\captionsetup[subfigure]{labelfont={sc}, skip=3pt}

% =========================================================================
%                     VARIABLES PARA LA PORTADA
% =========================================================================
\newcommand{\Universidad}{Universidad de Salamanca}
\newcommand{\Facultad}{Facultad de Ciencias}
\newcommand{\Titulacion}{Grado en Física}
\newcommand{\Asignatura}{Laboratorio de Mecánica y Ondas}
\newcommand{\NumPractica}{Práctica 8}
\newcommand{\TituloPractica}{Efecto Doppler\\con ultrasonidos}
\newcommand{\tipotrabajo}{Informe de laboratorio}
\newcommand{\miNombre}{David Gómez García \\ Daniel Martín Guerrero}
\newcommand{\miTutor}{Mario Amado \\ Ivan de Martino \\ Manuel Pino}
\newcommand{\miUbicacion}{Salamanca}

% =========================================================================
%                     DATOS DEL ARTÍCULO (cabecera)
% =========================================================================
\title{Efecto Doppler con ultrasonidos}
\author{David Gómez García}
\author{Daniel Martín Guerrero}
\address{Grado en Física, Facultad de Ciencias, Universidad de Salamanca,
         Salamanca, España}
\date{}
\keywords{efecto Doppler, ultrasonidos, velocidad del sonido,
          incertidumbre de medida, ajuste por mínimos cuadrados}



% =========================================================================
%                                DOCUMENTO
% =========================================================================
\begin{document}

% --- PORTADA ----------------------------------------------------------------
\pagenumbering{alph}
%%%%%%%%%%%%%%%%%%%%%%%%%%%%%%%%%%%%%%%%%%%%%%%%%%%%%%%%%%%%%%%%%%%%%%%%
% Portada — estilo artículo científico, motivo: frentes de onda Doppler
% Requiere (definidos en main.tex): colores azulprof, azulacero, azulclaro,
% grispizarra, grisclaro; y las macros \Universidad, \Facultad, \Titulacion,
% \Asignatura, \NumPractica, \TituloPractica, \tipotrabajo, \miNombre,
% \miTutor, \miUbicacion.
%%%%%%%%%%%%%%%%%%%%%%%%%%%%%%%%%%%%%%%%%%%%%%%%%%%%%%%%%%%%%%%%%%%%%%%%

\begin{titlepage}
\thispagestyle{empty}
\pagecolor{white}

\begin{tikzpicture}[remember picture, overlay,
    pg/.style={x=1cm, y=1cm}]
  \coordinate (O) at (current page.south west);

  % ----------------------------------------------------------------
  % Motivo gráfico: frentes de onda de una fuente que se desplaza hacia
  % la derecha (compresión a la derecha, dilatación a la izquierda).
  % ----------------------------------------------------------------
  \begin{scope}
    \clip (O) rectangle ($(O)+(21cm,24.6cm)$);
    \foreach \k in {0,...,15}{
      \pgfmathsetmacro{\cx}{8.2+0.36*\k}
      \pgfmathsetmacro{\rr}{8.4-0.60*\k}
      \pgfmathsetmacro{\pc}{14+5.2*\k}
      \pgfmathsetmacro{\lw}{0.25+0.035*\k}
      \draw[azulprof!\pc!azulclaro, line width=\lw pt]
        ($(O)+(\cx cm,17.2cm)$) circle (\rr cm);
    }
    \fill[azulprof] ($(O)+(8.2cm+0.36cm*15,17.2cm)$) circle (2.2pt);
  \end{scope}

  % ----------------------------------------------------------------
  % Señal con desplazamiento Doppler (chirp): frecuencia creciente
  % ----------------------------------------------------------------
  \begin{scope}[shift={(O)}]
    \draw[azulacero, line width=0.9pt, domain=2.2:18.8, samples=420, smooth]
      plot ({\x}, {12.95 + 0.33*sin(deg((\x-2.2)*(2.2+0.58*(\x-2.2))))});
  \end{scope}

  % ----------------------------------------------------------------
  % Cabecera: logotipos y filete
  % ----------------------------------------------------------------
  \node[anchor=north west, inner sep=0] at ($(O)+(2.2cm,27.9cm)$)
    {\includegraphics[height=1.42cm]{Logo USAL.png}};
  \node[anchor=north east, inner sep=0] at ($(O)+(18.8cm,27.9cm)$)
    {\includegraphics[height=1.42cm]{Logo Facultad de Ciencias.png}};
  \draw[grispizarra!60, line width=0.4pt]
    ($(O)+(2.2cm,26.0cm)$) -- ($(O)+(18.8cm,26.0cm)$);
  \node[anchor=north west, inner sep=0, text=grispizarra, font=\scshape\small]
    at ($(O)+(2.2cm,25.75cm)$)
    {\textls[80]{\Universidad\ · \Facultad}};
  \node[anchor=north east, inner sep=0, text=grispizarra, font=\scshape\small]
    at ($(O)+(18.8cm,25.75cm)$)
    {\textls[80]{\Titulacion}};

  % ----------------------------------------------------------------
  % Bloque del título
  % ----------------------------------------------------------------
  \node[anchor=south west, inner sep=0, text=azulacero, font=\scshape\large,
        text width=16.6cm]
    at ($(O)+(2.2cm,11.2cm)$)
    {\textls[120]{\Asignatura\ ·\ \NumPractica}};

  \node[anchor=north west, inner sep=0, text=azulprof, text width=16.6cm,
        align=left]
    at ($(O)+(2.2cm,11.0cm)$)
    {\fontsize{35}{39}\selectfont\bfseries\TituloPractica};

  \draw[azulacero, line width=1.4pt]
    ($(O)+(2.2cm,7.35cm)$) -- ($(O)+(4.7cm,7.35cm)$);

  \node[anchor=north west, inner sep=0, text=grispizarra, text width=16.6cm,
        align=left, font=\itshape\large]
    at ($(O)+(2.2cm,7.05cm)$)
    {$\displaystyle f_{\mathrm{R}} \;=\; f_{\mathrm{E}}\,
      \frac{v_{\mathrm{s}}-v_{\mathrm{R}}}{v_{\mathrm{s}}-v_{\mathrm{E}}}$};

  % ----------------------------------------------------------------
  % Pie: tipo de trabajo, autores, tutores, lugar y fecha
  % ----------------------------------------------------------------
  \draw[grispizarra!60, line width=0.4pt]
    ($(O)+(2.2cm,5.1cm)$) -- ($(O)+(18.8cm,5.1cm)$);

  \node[anchor=north west, inner sep=0, text width=4.6cm, align=left]
    at ($(O)+(2.2cm,4.75cm)$)
    {{\scshape\small\color{azulacero}\textls[80]{Documento}}\\[3pt]
     {\large\bfseries\color{azulprof}\tipotrabajo}};

  \node[anchor=north west, inner sep=0, text width=5.6cm, align=left]
    at ($(O)+(7.2cm,4.75cm)$)
    {{\scshape\small\color{azulacero}\textls[80]{Autores}}\\[3pt]
     {\large\color{black!85}\miNombre}};

  \node[anchor=north west, inner sep=0, text width=5cm, align=left]
    at ($(O)+(13.8cm,4.75cm)$)
    {{\scshape\small\color{azulacero}\textls[80]{Tutores}}\\[3pt]
     {\large\color{black!85}\miTutor}};

  \node[anchor=south east, inner sep=0, text=grispizarra, font=\scshape\small]
    at ($(O)+(18.8cm,1.45cm)$)
    {\textls[100]{\miUbicacion, \number\year}};
  \fill[azulprof] ($(O)+(0,0)$) rectangle ($(O)+(21cm,0.45cm)$);
  \fill[azulacero] ($(O)+(0,0)$) rectangle ($(O)+(7cm,0.45cm)$);
  \fill[azulclaro] ($(O)+(0,0)$) rectangle ($(O)+(3cm,0.45cm)$);
\end{tikzpicture}

\end{titlepage}

\clearpage
\pagenumbering{arabic}
\setcounter{page}{1}

\begin{abstract}
Se estudia el efecto Doppler acústico mediante ultrasonidos de frecuencia fija
emitidos y recibidos sobre un carril, con emisor o receptor en movimiento
uniforme. A partir de la relación entre las frecuencias emitida y recibida se
determina la velocidad de propagación del sonido en el aire por regresión
lineal y se contrasta con el valor teórico deducido de la temperatura
ambiente; se estima, además, la frecuencia de emisión en las dos
configuraciones posibles y se discute si el término cuadrático en $v/\vs$ es
resoluble con el equipo disponible. Todas las magnitudes se expresan con su
incertidumbre típica y ampliada al \SI{95}{\percent}, siguiendo la metodología de la
\emph{Guide to the expression of uncertainty in measurement}.
\end{abstract}

\maketitle
\tableofcontents

% =========================================================================
\section{Introducción}
% =========================================================================
La experiencia y su posterior análisis recogidos en este informe presentan
como temática nuclear un fenómeno físico conocido como efecto Doppler. Fue
propuesto por Christian Doppler en 1842 para explicar el color de las
estrellas dobles~\cite{DOPPLER-1842} y verificado para el sonido pocos años
después por Buys Ballot, quien hizo sonar una nota constante desde un
ferrocarril en marcha y comprobó con oído entrenado que el tono percibido
variaba según el tren se acercase o se alejase~\cite{BUYSBALLOT-1845}. Desde
entonces el fenómeno se ha convertido en una herramienta de medida
omnipresente: la velocimetría por ultrasonidos, el radar meteorológico y de
tráfico o el corrimiento al rojo de las galaxias descansan sobre el mismo
principio.

En esta práctica se reproduce el efecto en condiciones controladas. Un
emisor y un receptor de ultrasonidos se desplazan uno respecto del otro a
velocidad constante y conocida, de modo que la diferencia entre la frecuencia
emitida y la recibida puede medirse, compararse con la predicción teórica y
utilizarse para inferir la velocidad del sonido en el aire. El informe se
organiza como sigue: la \cref{sec:fundamentos} desarrolla el fundamento
teórico; la \cref{sec:procedimiento} describe el procedimiento y establece el
método de tratamiento de incertidumbres y de ajuste; la
\cref{sec:analisis} recoge el análisis de los datos y la \cref{sec:discusion}
su discusión y las conclusiones.

\subsection{Objetivos}\label{sec:objetivos}
Los objetivos de la práctica son los siguientes.
\begin{itemize}
    \item Hallar la velocidad con la que se propagan ondas ultrasonidos en un medio como el aire.
    \item Evidenciar y calcular las diferencias entre la frecuencia emitida y recibida ante varias situaciones dinámicas del emisor y receptor.
\end{itemize}
El primero se aborda por vía indirecta: la velocidad del sonido aparece como
pendiente de una recta cuando se representa convenientemente la velocidad
relativa frente a una función de las frecuencias (\cref{sec:casos}), y el
resultado se confronta con el valor teórico que predice la temperatura del
laboratorio. El segundo exige medir $\fE$, $\fR$ y la velocidad relativa
$v$ en cuatro situaciones dinámicas (emisor móvil y receptor en reposo, y
viceversa, en ambos sentidos de la marcha) y comprobar que se ajustan a la
ley de Doppler dentro de las incertidumbres experimentales.

\subsection{Material empleado}\label{sec:material}
Para el correcto desarrollo del montaje experimental, se ha hecho uso del siguiente instrumental:
\begin{itemize}
    \item Unidad de emisión y detección de ultrasonidos.
    \item Carril y coche autopropulsado.
    \item Emisor, receptor y soportes.
    \item Barrera fotoeléctrica y temporizador.
    \item Ordenador y unidad de medida Cobra-3.
\end{itemize}

\subsection{Montaje experimental}\label{sec:montaje}
El montaje (\cref{fig:montaje}) se articula en torno a un carril rectilíneo
sobre el que circula el coche autopropulsado, cuya velocidad se fija mediante
un selector mecánico y cuyo sentido de marcha se invierte con un interruptor.
El emisor de ultrasonidos y el receptor se sujetan a soportes verticales
enfrentados y alineados con el eje del carril; uno de ellos permanece fijo y
el otro se monta sobre el coche, de manera que la dirección que une emisor y
receptor coincide con la de movimiento, como exige el modelo
unidimensional de la \cref{sec:fundamentos}. La barrera fotoeléctrica se
coloca aproximadamente a la mitad del recorrido para que el coche la atraviese
a velocidad lo más uniforme posible. La unidad Cobra-3 se conecta al ordenador
y, a través del programa de medida, cumple dos funciones: con el modo
«Temporizador» determina la velocidad del coche a partir del tiempo de
oscuración de la barrera, y con el modo «Contador» mide la frecuencia de la
señal recibida~\cite{GUION-P8}.

\begin{figure}[ht]
\centering
\begin{tikzpicture}[>=Stealth, x=1cm, y=1cm, font=\small]
  % --- carril ---
  \fill[grisclaro] (0,0) rectangle (13.2,0.26);
  \draw[grispizarra, line width=0.6pt] (0,0) rectangle (13.2,0.26);
  \foreach \x in {0.6,1.2,...,12.6}{\draw[grispizarra!50, line width=0.3pt] (\x,0)--(\x,0.26);}
  % --- coche ---
  \begin{scope}[shift={(2.2,0.26)}]
    \draw[azulprof, fill=azulclaro!55, line width=0.7pt, rounded corners=1.5pt]
      (0,0.16) rectangle (2.5,0.78);
    \draw[azulprof, fill=white, line width=0.7pt] (0.5,0.16) circle (0.16);
    \draw[azulprof, fill=white, line width=0.7pt] (2.0,0.16) circle (0.16);
    % soporte y emisor
    \draw[grispizarra, line width=0.9pt] (1.25,0.78) -- (1.25,1.35);
    \draw[azulprof, fill=azulacero!85, line width=0.7pt]
      (1.25,1.35) -- ++(0.07,0.07) -- ++(0,0.46) -- ++(-0.14,0) -- ++(0,-0.46) -- cycle;
    \fill[azulprof] (1.31,1.60) circle (0.05);
    \node[azulprof, anchor=south] at (1.25,2.0) {emisor $(\mathrm{E})$};
    % velocidad
    \draw[->, line width=0.9pt, azulprof] (-0.05,1.1) -- (-0.85,1.1);
    \node[azulprof, anchor=east] at (-0.9,1.1) {$\vE$};
  \end{scope}
  % --- ondas emitidas ---
  \foreach \r/\o in {0.45/90, 0.85/70, 1.25/50, 1.65/35}{
    \draw[azulacero, opacity=0.35+\o/200, line width=0.7pt]
      (3.55,1.86) ++(-32:\r) arc[start angle=-32, end angle=32, radius=\r];
  }
  % --- receptor fijo ---
  \begin{scope}[shift={(11.0,0.26)}]
    \draw[grispizarra, line width=0.9pt] (0,0) -- (0,1.35);
    \draw[azulprof, fill=azulacero!85, line width=0.7pt]
      (0,1.35) -- ++(-0.07,0.07) -- ++(0,0.46) -- ++(0.14,0) -- ++(0,-0.46) -- cycle;
    \node[azulprof, anchor=south] at (0,2.0) {receptor $(\mathrm{R})$};
    \draw[grispizarra, fill=grisclaro, line width=0.6pt] (-0.4,0) rectangle (0.4,0.12);
  \end{scope}
  % --- barrera fotoeléctrica ---
  \begin{scope}[shift={(7.0,0.26)}]
    \draw[grispizarra, line width=1pt] (-0.45,0) -- (-0.45,1.05);
    \draw[grispizarra, line width=1pt] (0.45,0) -- (0.45,1.05);
    \fill[azulprof] (-0.45,0.95) circle (0.07);
    \fill[azulprof] (0.45,0.95) circle (0.07);
    \draw[azulacero, densely dashed, line width=0.6pt] (-0.45,0.95) -- (0.45,0.95);
    \node[azulprof, align=center, anchor=south] at (0,1.15)
      {barrera\\fotoeléctrica};
  \end{scope}
  % --- acotación de la distancia de propagación ---
  \draw[|<->|, grispizarra, line width=0.5pt] (3.7,-0.55) -- (11.0,-0.55)
    node[midway, fill=white, inner sep=1.5pt, text=black!80] {línea de propagación};
  % --- interfaz y ordenador ---
  \draw[azulprof, fill=azulclaro!55, line width=0.7pt, rounded corners=2pt]
    (5.6,-2.35) rectangle (8.0,-1.45);
  \node[azulprof, align=center] at (6.8,-1.9) {unidad\\Cobra-3};
  \draw[azulprof, fill=white, line width=0.7pt, rounded corners=2pt]
    (9.5,-2.45) rectangle (12.6,-1.45);
  \node[azulprof, align=center] at (11.05,-1.95) {ordenador\\PHYWE Measure};
  \draw[->, line width=0.8pt, azulprof] (8.0,-1.9) -- (9.5,-1.95);
  % --- conexiones (cables) ---
  \draw[grispizarra, densely dotted, line width=0.8pt]
    (7.0,0.26) -- (7.0,-1.45);
  \draw[grispizarra, densely dotted, line width=0.8pt]
    (11.0,0.26) .. controls (11.0,-0.9) and (8.4,-0.9) .. (8.0,-1.7);
  \draw[grispizarra, densely dotted, line width=0.8pt]
    (3.45,0.26) .. controls (3.45,-1.0) and (5.4,-1.0) .. (5.6,-1.8);
\end{tikzpicture}
\caption{Esquema del montaje experimental. El emisor, solidario al coche
autopropulsado, se mueve con velocidad $\vE$ a lo largo del carril; el
receptor permanece fijo (en la segunda configuración se intercambian sus
posiciones). La barrera fotoeléctrica, situada hacia la mitad del recorrido,
permite medir la velocidad del coche, y la unidad Cobra-3 adquiere tanto esa
velocidad como la frecuencia recibida y las transfiere al ordenador. Las
líneas de puntos representan las conexiones eléctricas.}
\label{fig:montaje}
\end{figure}

% =========================================================================
\section{Fundamentos teóricos}\label{sec:fundamentos}
% =========================================================================

\subsection{Descripción del fenómeno}
\begin{definicion}[Efecto Doppler]\label{def:doppler}
El efecto Doppler es un fenómeno físico que describe el cambio en la
frecuencia percibida de una onda cuando la fuente emisora de la onda y el
observador están en movimiento relativo entre sí.
\end{definicion}

Este cambio de frecuencia no se debe a una modificación real de la frecuencia
de la onda emitida, sino a un cambio en la forma en que el observador percibe
la onda debido al movimiento. Su causa radica en la compresión o dilatación
de las ondas debida al movimiento relativo entre la fuente y el observador.
Si la fuente se mueve hacia el observador, las ondas se comprimen, lo que
resulta en una longitud de onda más corta y una frecuencia percibida más
alta; en las ondas de sonido se aprecia una frecuencia más aguda. Por el
contrario, si la fuente se aleja, las ondas se dilatan, la longitud de onda
aparente es mayor y la frecuencia menor: el sonido se hace más grave. El
efecto se aplica a cualquier tipo de onda, incluyendo ondas sonoras, ondas de
luz y ondas de agua~\cite{FV6-ECJ}. Conviene advertir que, en el caso del
sonido, el movimiento debe referirse al medio de propagación: no es lo mismo
que se mueva la fuente que que lo haga el observador, aunque, como se verá,
ambas situaciones coinciden a primer orden en $v/\vs$.

La \cref{fig:frentes} ilustra el mecanismo para una fuente $\mathrm{S}$ que se
desplaza con velocidad $\vE>0$ hacia la derecha. En cada instante $t_k$ la
fuente se encuentra en $x_k=\vE t_k$ y emite un frente que, transcurrido un
tiempo $t-t_k$, ha alcanzado el radio $r_k=\vs\,(t-t_k)$. Los frentes se
amontonan por delante de la fuente y se separan por detrás; un receptor
situado delante, $\mathrm{R}_2$, registra por tanto una señal más rápida que
la emitida, y uno situado detrás, $\mathrm{R}_1$, una más lenta.

\begin{figure}[ht]
\centering
\begin{tikzpicture}[>=Stealth, scale=0.85, every node/.style={scale=0.9}]
  \definecolor{cS}{HTML}{1B365D}
  \definecolor{cA}{HTML}{7A8C9E}
  \definecolor{cB}{HTML}{3B6EA8}
  % ---------- Panel izquierdo: frentes espaciales ----------
  \begin{scope}
    \node[anchor=west, text=grispizarra] at (-5.6,3.9) {\small\scshape (a)\; frentes de onda};
    \draw[->, line width=0.7pt, grispizarra] (-5.5,0) -- (4.6,0) node[right, text=black] {$x$};
    \def\rA{3.8} \def\cA{0}
    \def\rB{2.8} \def\cB{0.7}
    \def\rC{1.8} \def\cC{1.4}
    \def\rD{0.8} \def\cD{2.1}
    \foreach \c/\r/\o in {\cA/\rA/45,\cB/\rB/60,\cC/\rC/75,\cD/\rD/90}{
      \draw[cB!\o!cA, line width=0.9pt] (\c,0) circle (\r);}
    \foreach \c/\n in {\cA/0,\cB/1,\cC/2,\cD/3}{
      \fill[cS] (\c,0) circle (2.3pt) node[below=3pt, text=cS] {$x_{\n}$};}
    \node[text=cS, above=4pt] at (\cD,0) {$\mathrm{S}$};
    \fill[cB] (-4.7,0) circle (3pt) node[below=4pt, text=cB] {$\mathrm{R}_1$};
    \fill[cB] (3.55,0) circle (3pt) node[below=4pt, text=cB] {$\mathrm{R}_2$};
    \draw[->, dashed, grispizarra] (\cA,0) -- ++(110:\rA)
      node[midway, left=2pt, fill=white, inner sep=1pt, text=black!80] {$r_0$};
    \draw[->, dashed, grispizarra] (\cB,0) -- ++(80:\rB)
      node[midway, right=1pt, fill=white, inner sep=1pt, text=black!80] {$r_1$};
    \draw[->, dashed, grispizarra] (\cC,0) -- ++(50:\rC)
      node[midway, right=1pt, fill=white, inner sep=1pt, text=black!80] {$r_2$};
    \draw[->, dashed, grispizarra] (\cD,0) -- ++(20:\rD)
      node[midway, above=1pt, fill=white, inner sep=1pt, text=black!80] {$r_3$};
    \draw[->, line width=0.9pt, cS] (\cD-0.6,-1.7) -- ++(1.2,0)
      node[midway, below=1pt] {$\vE$};
    \node[anchor=west, align=left, text=black!80] at (-5.7,2.45)
      {$r_k=\vs\,(t-t_k)$};
    \node[anchor=west, align=left, text=black!80] at (-5.7,-2.45)
      {$x_k=\vE\,t_k$};
  \end{scope}
  % ---------- Panel derecho: señales ----------
  \begin{scope}[shift={(6.6,0)}]
    \def\w{5.5} \def\ha{1.1}
    \node[anchor=west, text=grispizarra] at (0,3.9) {\small\scshape (b)\; señales temporales};
    \foreach \y/\col/\lab/\f/\st in {2.55/cS/{\mathrm{S}}/1.5/0,
                                     0/cB/{\mathrm{R}_1}/0.9/1.8,
                                     -2.55/cB/{\mathrm{R}_2}/2.5/0.8}{
      \begin{scope}[shift={(0,\y)}]
        \draw[grisclaro, line width=0.4pt, step=0.5] (0,-\ha) grid (\w,\ha);
        \draw[->, line width=0.7pt, grispizarra] (0,0) -- (\w+0.45,0);
        \draw[->, line width=0.7pt, grispizarra] (0,-\ha-0.2) -- (0,\ha+0.2);
        \draw[\col, line width=0.9pt] (0,0) -- (\st,0);
        \draw[\col, line width=0.9pt, domain=\st:\w, samples=160]
          plot (\x, {0.85*sin((\x-\st)*180*\f)});
        \node[\col, fill=white, inner sep=2.5pt, draw=\col!60, line width=0.4pt]
          at (\w-0.45,\ha-0.2) {$\lab$};
      \end{scope}}
    \node[text=black!80, anchor=north] at (\w/2,-3.9) {tiempo $t$};
  \end{scope}
\end{tikzpicture}
\caption{Efecto Doppler generado por una fuente $\mathrm{S}$ que se desplaza
con velocidad $\vE$. (a)~Frentes de onda emitidos en distintos instantes:
compresión en el sentido del movimiento y dilatación en el opuesto.
(b)~Señal emitida frente a las percibidas por los receptores $\mathrm{R}_1$
(menor frecuencia) y $\mathrm{R}_2$ (mayor frecuencia), desfasadas por el
tiempo de propagación. Esquema cualitativo, sin escala.}
\label{fig:frentes}
\end{figure}

\subsection{Deducción en el sistema de referencia de la onda}\label{sec:deduccion}
Para cuantificar el efecto se adopta un modelo unidimensional con las
hipótesis siguientes.
\begin{hipotesis}\label{hip:modelo}
El medio de propagación está en reposo y emisor y receptor se mueven a lo
largo de la recta que los une con velocidades constantes $\vE$ y $\vR$
respecto al medio (positivas o negativas según el sentido del movimiento,
tomando como positivo el de propagación, del emisor al receptor). La onda se
propaga con velocidad constante $\vs>\vE,\,\vR$ respecto al medio.
\end{hipotesis}

Analicemos la emisión y la detección en el sistema de referencia de la onda
(\cref{fig:relativo}). En dicho sistema el emisor se mueve con velocidad
$\vE-\vs<0$ y el receptor con velocidad $\vR-\vs<0$. El emisor se desplaza una
distancia igual a una longitud de onda $\lambda$ en un intervalo $\TE$ tal que
\begin{equation}\label{eq:lambdaE}
  \lambda = \TE\,(\vs-\vE),
\end{equation}
y el receptor recorre una longitud de onda en un intervalo $\TR$ tal que
\begin{equation}\label{eq:lambdaR}
  \lambda = \TR\,(\vs-\vR).
\end{equation}
Los tiempos $\TE$ y $\TR$ son, respectivamente, los períodos de la onda emitida
y de la observada.

\begin{figure}[ht]
\centering
\begin{tikzpicture}[>=Stealth, scale=1.15]
  \draw[->, line width=0.5pt, grispizarra!60] (-2,0) -- (10,0)
    node[right, text=grispizarra] {$x$};
  \draw[azulacero, line width=0.9pt, densely dashed, domain=-1.5:1, samples=60]
    plot (\x,{sin(\x*180)});
  \draw[azulacero, line width=0.9pt, densely dashed, domain=7:9.5, samples=60]
    plot (\x,{sin(\x*180)});
  \draw[azulprof, line width=1.4pt, domain=1:7, samples=180]
    plot (\x,{sin(\x*180)});
  \foreach \x/\t in {1/emisor,7/receptor}{
    \draw[grispizarra!70, densely dotted, line width=0.8pt] (\x,-1.4) -- (\x,1.4);
    \node[draw=grispizarra!50, fill=white, inner sep=3.5pt, text=azulprof,
          font=\small\scshape] at (\x,1.75) {\t};}
  \draw[->, line width=0.9pt, azulprof] (1,-1.75) -- (0,-1.75)
    node[midway, below] {$\vE-\vs$};
  \draw[->, line width=0.9pt, azulprof] (7,-1.75) -- (6,-1.75)
    node[midway, below] {$\vR-\vs$};
  \draw[grispizarra!60, line width=0.4pt] (2.5,1.05) -- (2.5,1.55);
  \draw[grispizarra!60, line width=0.4pt] (4.5,1.05) -- (4.5,1.55);
  \draw[|<->|, line width=0.7pt, azulacero] (2.5,1.4) -- (4.5,1.4)
    node[midway, fill=white, inner sep=2pt, above=-1pt] {$\lambda$};
\end{tikzpicture}
\caption{Movimientos del emisor y del receptor en el sistema de referencia
de la onda, para $\vs>\vE,\vR$. Se toma como sentido positivo de propagación
el de izquierda a derecha; ambas velocidades relativas son negativas.}
\label{fig:relativo}
\end{figure}

Dividiendo ambas expresiones se obtiene la relación entre los períodos y,
equivalentemente, entre las frecuencias, $f=1/T$:
\begin{equation}\label{eq:cociente}
  \frac{\TR}{\TE}=\frac{\vs-\vE}{\vs-\vR}
  \qquad\Longleftrightarrow\qquad
  \frac{\fR}{\fE}=\frac{\vs-\vR}{\vs-\vE}.
\end{equation}
Resulta cómodo reescribir la relación en términos de la velocidad relativa
entre emisor y receptor,
\begin{equation}\label{eq:vrel}
  v \equiv \vE-\vR,
\end{equation}
que es positiva cuando emisor y receptor se acercan y negativa cuando se
alejan, con independencia de los sentidos individuales de movimiento. Como
$\vs-\vR=(\vs-\vE)+v$, la \cref{eq:cociente} adopta la forma
\begin{equation}\label{eq:doppler}
  \fR=\fE\left(1+\frac{v}{\vs-\vE}\right).
\end{equation}

\begin{proposicion}[Ley de Doppler para ondas sonoras]\label{prop:doppler}
Bajo la \cref{hip:modelo}, las frecuencias emitida y recibida están
relacionadas por la \cref{eq:doppler}. En particular, $\fR>\fE$ si emisor y
receptor se acercan ($v>0$) y $\fR<\fE$ si se alejan ($v<0$).
\end{proposicion}

\subsection{Casos particulares y linealización}\label{sec:casos}
Los dos casos que se estudian experimentalmente corresponden a que se mueva
solo el emisor o solo el receptor.

\medskip\noindent\textbf{Emisor en movimiento y receptor en reposo ($\vE\neq0,\ \vR=0$).}
Entonces $v=\vE$ y la \cref{eq:doppler} se reduce a
\begin{equation}\label{eq:emisor}
  \fR=\frac{\fE}{1-v/\vs}.
\end{equation}
Despejando la velocidad relativa resulta
\begin{equation}\label{eq:vsdespeje}
  v=\vs\left(1-\frac{\fE}{\fR}\right),
\end{equation}
una relación \emph{lineal} entre $v$ y $1-\fE/\fR$ cuya pendiente es la
velocidad de propagación de la onda. Esta expresión permite, por tanto,
obtener $\vs$ midiendo la velocidad relativa y las frecuencias emitida y
observada.

\medskip\noindent\textbf{Emisor en reposo y receptor en movimiento ($\vE=0,\ \vR\neq0$).}
Ahora $v=-\vR$ y
\begin{equation}\label{eq:receptor}
  \fR=\fE\left(1+\frac{v}{\vs}\right).
\end{equation}
Con la notación $x\equiv v/\vs$ y la identidad
$1+x=\dfrac{1}{1-x}-\dfrac{x^{2}}{1-x}$, la \cref{eq:receptor} se escribe
\begin{equation}\label{eq:receptor2}
  \fR=\frac{\fE}{1-v/\vs}-\fE\,\frac{(v/\vs)^{2}}{1-v/\vs}
  \;\approx\;\frac{\fE}{1-v/\vs},\qquad |v/\vs|\ll1.
\end{equation}

\begin{observacion}[Equivalencia a primer orden]\label{obs:primerorden}
A primer orden en $|v/\vs|$ las situaciones de emisor móvil y de receptor
móvil son indistinguibles: las \cref{eq:emisor,eq:receptor} difieren
únicamente en el término cuadrático. El efecto Doppler acústico es, pues,
asimétrico en sentido estricto (importa qué elemento se mueve respecto al
medio), aunque la asimetría solo es apreciable a velocidades comparables a la
del sonido.
\end{observacion}

\begin{proposicion}[Término cuadrático]\label{prop:cuadratico}
Sea $\fR$ la frecuencia recibida con el receptor en movimiento y el emisor en
reposo. La diferencia respecto a la predicción de emisor móvil vale
\begin{equation}\label{eq:difcuadratica}
  \fR-\frac{\fE}{1-v/\vs}=-\fE\,\frac{(v/\vs)^{2}}{1-v/\vs}\;\approx\;-\fE\Bigl(\frac{v}{\vs}\Bigr)^{2}.
\end{equation}
\end{proposicion}
\begin{proof}
Es consecuencia inmediata de la identidad algebraica anterior,
$(1+x)-1/(1-x)=-x^{2}/(1-x)$, con $x=v/\vs$.
\end{proof}

La \cref{fig:cuadratico} compara ambas leyes. La diferencia relativa entre
ellas crece como $x^{2}$, frente a la contribución lineal $x$, por lo que su
observación exige medir con una precisión relativa en la frecuencia
comparable a $|v/\vs|^{2}$. Esta cuestión se retoma en el análisis de datos.

\begin{figure}[ht]
\centering
\pgfplotsset{
  doppler/.style={
    width=0.40\textwidth, height=5.4cm, tick label style={font=\scriptsize},
    label style={font=\small}, axis lines=left, axis line style={grispizarra},
    grid=major, grid style={grisclaro, line width=0.4pt}, every tick/.style={grispizarra},
    legend style={font=\scriptsize, draw=none, fill=white, fill opacity=0.85,
                  text opacity=1, at={(0.03,0.97)}, anchor=north west},
    xlabel={$x=v/\vs$}}}
\begin{tikzpicture}
  \begin{axis}[doppler, ylabel={$\fR/\fE$}, xmin=-0.3, xmax=0.3,
               ymin=0.6, ymax=1.5, xtick={-0.3,-0.15,0,0.15,0.3}]
    \addplot[azulprof, line width=1pt, domain=-0.3:0.3, samples=100] {1/(1-x)};
    \addplot[azulacero, line width=1pt, dashed, domain=-0.3:0.3, samples=2] {1+x};
    \addplot[grispizarra, line width=0.6pt, densely dotted, domain=-0.3:0.3, samples=2]
      {1+x}; % guía de primer orden
    \legend{emisor móvil \eqref{eq:emisor},receptor móvil \eqref{eq:receptor}}
  \end{axis}
\end{tikzpicture}\hfill
\begin{tikzpicture}
  \begin{axis}[doppler, ylabel={$\bigl(\fR-\fE/(1-x)\bigr)/\fE$}, xmin=-0.3, xmax=0.3,
               ymin=-0.14, ymax=0.01, xtick={-0.3,-0.15,0,0.15,0.3},
               scaled y ticks=false, yticklabel style={/pgf/number format/fixed,
               /pgf/number format/precision=2}]
    \addplot[azulprof, line width=1pt, domain=-0.3:0.3, samples=100]
      {-x^2/(1-x)};
    \addplot[grispizarra, line width=0.7pt, densely dashed, domain=-0.3:0.3, samples=2] {0};
  \end{axis}
\end{tikzpicture}
\caption{Comparación de las leyes de emisor móvil y de receptor móvil.
(Izquierda)~Cociente $\fR/\fE$ en función de $x=v/\vs$; la línea punteada es
el desarrollo a primer orden, común a ambos casos. (Derecha)~Diferencia
relativa entre ambas leyes \eqref{eq:difcuadratica}, nula para el emisor
móvil por construcción. El rango de $x$ se ha ampliado con fines ilustrativos;
en el laboratorio $|x|\ll1$.}
\label{fig:cuadratico}
\end{figure}

% =========================================================================
\section{Procedimiento experimental}\label{sec:procedimiento}
% =========================================================================

\subsection{Protocolo de medida}\label{sec:protocolo}
Tras encender la fuente de alimentación, el ordenador y la unidad de medida,
se ejecuta el programa PHYWE Measure y se selecciona el tipo de medida
«Temporizador/Contador», sin modificar el resto de opciones. El modo
«Temporizador», en automático, activa la barrera fotoeléctrica y determina la
velocidad del coche en el instante en que la cruza; el modo «Contador», en
manual, registra la frecuencia recibida cuando se pulsa la barra
espaciadora. Cuando emisor o receptor van sobre el coche, la frecuencia se
mide en el instante en que este atraviesa la barrera, de modo que velocidad y
frecuencia corresponden al mismo punto del recorrido~\cite{GUION-P8}.

\textbf{Emisor en movimiento y receptor en reposo.} Se determina en primer
lugar la frecuencia de la señal en reposo, $\fE$, con al menos seis medidas.
Se fija después una velocidad del coche con el selector mecánico y se miden
seis veces tanto esa velocidad como la frecuencia observada, para cada
sentido de la marcha, ajustando la amplitud de emisión para que la
recepción sea adecuada. El proceso se repite para tres posiciones distintas
del selector de velocidad.

\textbf{Emisor en reposo y receptor en movimiento.} Se intercambian los
soportes verticales de emisor y receptor, de modo que este último quede
sobre el coche, y se sigue exactamente el mismo protocolo.

\subsection{Velocidad del sonido esperada}\label{sec:vsesperada}
En primer lugar se ha comenzado la experiencia midiendo el valor de la
temperatura del aire, $T=\pend{\ldots}\,\si{\celsius}$, a fin de determinar el
valor teórico o valor esperado de la velocidad del sonido en el aire,
$\vst$, mediante
\begin{equation}\label{eq:vs}
  \vst = \SI{331,3}{\meter\per\second}\,\sqrt{1+\frac{T}{\SI{273,15}{}}},
\end{equation}
donde $T$ es la temperatura en grados Celsius. Este valor se emplea como
referencia con la que contrastar la velocidad del sonido obtenida del ajuste
y, en el análisis posterior, como dato en las abscisas $v/\vst$.

\subsection{Tratamiento de incertidumbres}\label{sec:incertidumbres}
El tratamiento de incertidumbres sigue las recomendaciones de la \emph{Guide to
the expression of uncertainty in measurement}~\cite{GUM}, en la presentación
de las notas del laboratorio~\cite{NUM-ARG}; para el vocabulario se sigue
el~\cite{VIM}. Se resumen a continuación las definiciones empleadas.

\begin{definicion}[Incertidumbre típica de una medida directa]\label{def:tipica}
Sean $x_1,\dots,x_n$ medidas independientes de una misma magnitud $x$. Se toma
como mejor estimación la media muestral y como incertidumbre típica la
desviación típica de la media,
\begin{equation}\label{eq:media}
  \mean{x}=\frac1n\sum_{i=1}^{n}x_i,\qquad
  s_x=\sqrt{\frac{1}{n-1}\sum_{i=1}^{n}(x_i-\mean{x})^{2}},\qquad
  u(\mean{x})=\frac{s_x}{\sqrt n},
\end{equation}
con $\nu=n-1$ grados de libertad.
\end{definicion}

\begin{definicion}[Límite de resolución]\label{def:resolucion}
Si el instrumento tiene una resolución $\Delta$, se supone que el valor
verdadero se distribuye uniformemente en un intervalo de anchura $\Delta$
centrado en la lectura, cuya varianza es $\Delta^{2}/12$. Combinada con la
incertidumbre estadística,
\begin{equation}\label{eq:resolucion}
  u(x_{\mathrm{b}})=\sqrt{u^{2}(\mean{x})+\frac{\Delta^{2}}{12}},
  \qquad
  \nu=(n-1)\,\frac{\bigl[u^{2}(\mean{x})+\Delta^{2}/12\bigr]^{2}}{u^{4}(\mean{x})}.
\end{equation}
En una lectura única se tiene $u=\Delta/\sqrt{12}$, con un número de grados de
libertad que se toma infinito.
\end{definicion}

\begin{proposicion}[Propagación de incertidumbres]\label{prop:propagacion}
Sea $y=f(x_1,\dots,x_N)$ una magnitud derivada de entradas no correlacionadas
con incertidumbres típicas $u(x_k)$. Su incertidumbre típica es
\begin{equation}\label{eq:propagacion}
  u(y)=\sqrt{\sum_{k=1}^{N}c_k^{2}\,u^{2}(x_k)},\qquad
  c_k=\frac{\partial f}{\partial x_k},
\end{equation}
y el número efectivo de grados de libertad, dado por la fórmula de
Welch--Satterthwaite, es
\begin{equation}\label{eq:welch}
  \nu_{\mathrm{eff}}=\frac{u^{4}(y)}{\sum_{k}c_k^{4}\,u^{4}(x_k)/\nu_k}.
\end{equation}
\end{proposicion}

\begin{definicion}[Incertidumbre ampliada]\label{def:ampliada}
La incertidumbre ampliada con un nivel de confianza del \SI{95}{\percent} es
\begin{equation}\label{eq:ampliada}
  \ext{y}=\Tcoef\,u(y),
\end{equation}
donde $\Tcoef$ es el factor de cobertura (valor de la $t$ de Student) para
$\nu$ grados de libertad (\cref{tab:student}). El resultado se expresa como
$y\pm\ext{y}$, indicando siempre el significado de la incertidumbre
citada y con no más de dos cifras significativas.
\end{definicion}

\begin{table}[ht]
\centering
\caption{Factor de cobertura $t_{95,\nu}$ para distintos grados de libertad $\nu$~\cite{NUM-ARG}.}
\label{tab:student}
\small
\begin{tabular}{@{}c *{14}{c}@{}}
\toprule
$\nu$ & 2 & 3 & 4 & 5 & 6 & 7 & 8 & 9 & 10 & 15 & 20 & 50 & 100 & $\infty$\\
\midrule
$t_{95,\nu}$ & \num{4,30} & \num{3,18} & \num{2,78} & \num{2,57} & \num{2,45} & \num{2,36} & \num{2,31} & \num{2,26} & \num{2,23} & \num{2,13} & \num{2,09} & \num{2,01} & \num{1,98} & \num{1,96}\\
\bottomrule
\end{tabular}
\end{table}

Con ello se calculan, en primer lugar, las incertidumbres asociadas a la
temperatura y a la velocidad esperada del sonido.

\begin{calculo}[Lectura de la temperatura]\label{calc:T}
Con una resolución del termómetro $\Delta=\SI{0,1}{\celsius}$ y una lectura
única, la \cref{eq:resolucion} da
\begin{equation}\label{eq:uT}
  u^{2}(T)=\frac{\Delta^{2}}{12}
  \quad\Longrightarrow\quad
  u(T)=\frac{\Delta}{\sqrt{12}}\approx\SI{0,03}{\celsius}.
\end{equation}
En consecuencia, el registro térmico al comienzo de la experiencia ha sido
$T=\pend{$t$}\pm\SI{0,03}{\celsius}$ (incertidumbre típica).
\end{calculo}

\begin{calculo}[Velocidad esperada del sonido]\label{calc:vs}
Determinar la incertidumbre asociada a la lectura de la temperatura es un paso
clave antes de poder hallar por propagación la incertidumbre propia de $\vst$.
Al depender $\vst$ solo de $T$, la \cref{eq:propagacion} se reduce a
\begin{equation}\label{eq:uvs0}
  u^{2}(\vst)=\left(\frac{\partial\vst}{\partial T}\right)^{\!2}u^{2}(T).
\end{equation}
La derivada parcial se obtiene a partir de la dependencia térmica de la
\cref{eq:vs}:
\begin{equation}\label{eq:dvsdT}
  \frac{\partial\vst}{\partial T}
  =\frac{331{,}3}{2\cdot273{,}15\,\sqrt{1+T/273{,}15}}
  =\frac{331{,}3}{546{,}3\,\sqrt{1+T/273{,}15}}
  \approx\frac{\SI{0,606}{\meter\per\second\per\kelvin}}{\sqrt{1+T/273{,}15}},
\end{equation}
que debe evaluarse en la temperatura medida. Sustituyendo en la
\cref{eq:uvs0} y extrayendo la raíz cuadrada,
\begin{equation}\label{eq:uvs}
  u(\vst)=\left|\frac{\partial\vst}{\partial T}\right|\frac{\Delta}{\sqrt{12}}
  =\left|\frac{\partial\vst}{\partial T}\right|\frac{\SI{0,1}{\kelvin}}{\sqrt{12}}
  \approx\frac{\SI{0,0175}{\meter\per\second}}{\sqrt{1+T/273{,}15}}
  =\pend{\ldots}\,\si{\meter\per\second}.
\end{equation}
Luego se espera que la velocidad del sonido en el aire para las condiciones
del laboratorio sea
\begin{equation}\label{eq:vsresultado}
  \vst=\pend{$\vst$}\pm\pend{$u(\vst)$}\ \si{\meter\per\second}
  \quad(\text{incertidumbre típica}).
\end{equation}
\end{calculo}

\begin{observacion}[Ámbito de validez]
La incertidumbre de la \cref{eq:uvs} recoge únicamente la procedente de la
lectura de $T$; no incluye la incertidumbre del propio modelo de la
\cref{eq:vs}, que supone aire seco, ni posibles gradientes de temperatura a lo
largo del carril. Por esta razón se espera que el contraste entre $\vst$ y la
velocidad ajustada se sitúe en el ámbito de la incertidumbre ampliada de esta
última.
\end{observacion}

\subsection{Estimación lineal por mínimos cuadrados}\label{sec:ajuste}
Tres de las cuestiones del análisis consisten en ajustar rectas a conjuntos de
datos. El modelo, y las expresiones para estimar los parámetros y sus
incertidumbres, son los siguientes~\cite{NUM-ARG}.

\begin{estimacion}[Ajuste lineal con incertidumbres comunes]\label{est:comun}
Dados $n$ pares $(x_i,y_i)$ con incertidumbre común en $y_i$ y despreciable en
$x_i$, el modelo $y=ax+b$ que minimiza $\sum_i(y_i-ax_i-b)^{2}$ tiene por
parámetros
\begin{equation}\label{eq:ajuste}
  a=\frac{\cov(x,y)}{\var(x)}=\frac{\mean{xy}-\mean{x}\,\mean{y}}{\mean{x^{2}}-\mean{x}^{2}},
  \qquad b=\mean{y}-a\,\mean{x},
\end{equation}
y sus incertidumbres típicas, con $\nu=n-2$ grados de libertad, son
\begin{equation}\label{eq:uajuste}
  s_\varepsilon=\sqrt{\frac{\sum_i(y_i-ax_i-b)^{2}}{n-2}},\qquad
  u(a)=\frac{s_\varepsilon}{\sqrt{(n-1)\var(x)}},\qquad
  u(b)=u(a)\sqrt{\mean{x^{2}}}.
\end{equation}
\end{estimacion}

\begin{estimacion}[Ajuste lineal con incertidumbres distintas]\label{est:pesos}
Si cada $y_i$ tiene incertidumbre $u_i$, se minimiza
$\chi^{2}=\sum_iw_i(y_i-ax_i-b)^{2}$ con pesos $w_i=1/u_i^{2}$. Las
expresiones~\eqref{eq:ajuste} siguen siendo válidas sustituyendo las medias
aritméticas por medias ponderadas $\langle\cdot\rangle$, y
\begin{equation}\label{eq:uajustepesos}
  u(a)=\frac{1}{\sqrt{\sum_iw_i}\,\sqrt{\langle x^{2}\rangle-\langle x\rangle^{2}}},\qquad
  u(b)=u(a)\sqrt{\langle x^{2}\rangle}.
\end{equation}
\end{estimacion}

Los parámetros se expresan como $a\pm\ext{a}$ y $b\pm\ext{b}$ con
$\ext{a}=\Tcoef\,u(a)$ y $\nu=n-2$. Si la incertidumbre de la abscisa no es
despreciable, se propaga al eje de ordenadas mediante la incertidumbre
efectiva $u_{\mathrm{ef},i}^{2}=u^{2}(y_i)+a^{2}u^{2}(x_i)$, que se itera hasta
la convergencia de $a$. El coeficiente de correlación
$r=\cov(x,y)/\sqrt{\var(x)\var(y)}$ y los residuos permiten juzgar la
aceptabilidad del modelo lineal.

Las tres rectas del análisis quedan recogidas en la \cref{tab:ajustes}; en
todas ellas, si el modelo es correcto, la ordenada en el origen debe ser
compatible con cero.

\begin{table}[ht]
\centering
\caption{Linealizaciones empleadas. El valor teórico del cociente $\fR/\fE$ en
cada caso procede de las \cref{eq:vsdespeje,eq:emisor,eq:receptor}.}
\label{tab:ajustes}
\small
\renewcommand{\arraystretch}{1.35}
\begin{tabular}{@{}llccc@{}}
\toprule
Configuración & Ordenada $y$ & Abscisa $x$ & Pendiente $a$ & Ordenada $b$\\
\midrule
Emisor móvil & $v$ & $1-\fE/\fR$ & $\vs$ & $0$\\
Emisor móvil & $\fR$ & $(1-v/\vst)^{-1}$ & $\fE$ & $0$\\
Receptor móvil & $\fR$ & $1+v/\vst$ & $\fE$ & $0$\\
\bottomrule
\end{tabular}
\end{table}

\begin{calculo}[Propagación a las variables del ajuste]\label{calc:abscisa}
Para el primer ajuste, la abscisa $x=1-\fE/\fR$ hereda las incertidumbres de las
frecuencias:
\begin{equation}\label{eq:ux}
  u^{2}(x)=\frac{u^{2}(\fE)}{\fR^{2}}+\frac{\fE^{2}\,u^{2}(\fR)}{\fR^{4}}.
\end{equation}
Para la diferencia $\Delta f=\fR-\fE\,(1-v/\vst)^{-1}$, empleada en la última
cuestión y con $x\equiv v/\vst$, resulta
\begin{equation}\label{eq:udf}
  u^{2}(\Delta f)=u^{2}(\fR)+\frac{u^{2}(\fE)}{(1-x)^{2}}
  +\frac{\fE^{2}}{\vst^{2}(1-x)^{4}}\Bigl[u^{2}(v)+x^{2}u^{2}(\vst)\Bigr].
\end{equation}
\end{calculo}

% =========================================================================
\section{Análisis de datos y resultados}\label{sec:analisis}
% =========================================================================
Siguiendo las indicaciones del guion~\cite{GUION-P8} y las recomendaciones de
las notas de incertidumbres~\cite{NUM-ARG}, esta sección reúne los datos
directos, los datos empleados en los ajustes y los resultados, con sus
incertidumbres típicas y ampliadas.

\subsection{Mediciones directas}
\Cref{tab:directas} muestra la plantilla de recogida de datos, donde se
indican explícitamente los límites de resolución de cada instrumento.

\begin{table}[ht]
\centering
\caption{Mediciones directas. Se indican la resolución $\Delta$ de cada
instrumento y las incertidumbres típicas asociadas.}
\label{tab:directas}
\small
\renewcommand{\arraystretch}{1.3}
\begin{tabular}{@{}lcccc@{}}
\toprule
Magnitud & $n$ & Valor & $u$ & $\Delta$\\
\midrule
Temperatura $T$ ($\si{\celsius}$) & 1 & \pendiente{ } & \SI{0,03}{} & \num{0,1}\\
Frecuencia emitida $\fE$ ($\si{\hertz}$) & $\geq6$ & \pendiente{ } & \pendiente{ } & \pendiente{ }\\
Velocidad $v$ ($\si{\meter\per\second}$) & 6 & \pendiente{ } & \pendiente{ } & \pendiente{ }\\
Frecuencia recibida $\fR$ ($\si{\hertz}$) & 6 & \pendiente{ } & \pendiente{ } & \pendiente{ }\\
\bottomrule
\end{tabular}
\end{table}

\subsection{Emisor en movimiento y receptor en reposo}
\textbf{Velocidad del sonido.} Se representa $v$ frente a $1-\fE/\fR$ y, según
la \cref{eq:vsdespeje}, se ajusta una recta cuya pendiente estima $\vs$ y se
da su intervalo de confianza al \SI{95}{\percent}, que se compara con $\vst$
(véase la \cref{eq:vsresultado}). \pendiente{Resultado del ajuste: $\vs=\ldots\pm\ldots$}

\textbf{Frecuencia de emisión.} Se representa $\fR$ frente a $(1-v/\vst)^{-1}$
y la pendiente del ajuste estima $\fE$, que se compara con el valor medido
directamente. \pendiente{Resultado del ajuste: $\fE=\ldots\pm\ldots$}

\subsection{Emisor en reposo y receptor en movimiento}
Se representa $\fR$ frente a $1+v/\vst$ y la pendiente del ajuste estima
$\fE$, que se compara con el valor medido directamente.
\pendiente{Resultado del ajuste: $\fE=\ldots\pm\ldots$}

\subsection{Análisis del término cuadrático}
Con los valores directamente medidos de $\fR$, $\fE$, $v$ y con $\vst$, se
calculan las diferencias $\Delta f$ y sus incertidumbres (\cref{eq:udf}) y se
representan frente a $v/\vst$ sin ajuste alguno. Según la
\cref{prop:cuadratico}, el término cuadrático solo es observable si
$|\Delta f|$ supera su incertidumbre ampliada $\ext{\Delta f}$.
\pendiente{Representación de $\Delta f$ y discusión}

% =========================================================================
\section{Discusión y conclusiones}\label{sec:discusion}
% =========================================================================
\begin{observacion}[Fuentes de incertidumbre]
Se han de valorar, entre otras, la uniformidad de la velocidad del coche en
las inmediaciones de la barrera, la sincronización entre la medida de $v$ y la
de $\fR$, la alineación de emisor y receptor con el eje del carril y las
reflexiones del ultrasonido en el entorno.
\end{observacion}

\begin{conclusion}\label{conc:final}
\pendiente{Conclusiones del informe: valor de $\vs$ y su compatibilidad con
$\vst$; frecuencias de emisión estimadas y su compatibilidad con las medidas
directas; observabilidad del término cuadrático.}
\end{conclusion}

% =========================================================================
%                              BIBLIOGRAFÍA
% =========================================================================
\printbibliography[title={Referencias}]

\end{document}
